\documentclass[11pt, a4paper]{article}

% ==============================================================================
% CONFIGURAÇÕES DE PÁGINA E IDIOMA
% ==============================================================================
\usepackage[a4paper, margin=2.5cm]{geometry}
\usepackage{fontspec}
\usepackage[brazilian, english, bidi=basic, provide=*]{babel}
\babelprovide[import, onchar=ids fonts]{brazilian}
\babelfont{rm}{Noto Sans}

% ==============================================================================
% PACOTES ESSENCIAIS PARA ESCRITA CIENTÍFICA
% ==============================================================================
% Matemática Avançada
\usepackage{amsmath}
\usepackage{amssymb}
\usepackage{amsfonts}
\usepackage{amsthm}
\usepackage{mathtools}

% Tabelas de Alta Qualidade
\usepackage{booktabs}
\usepackage{array}
\usepackage{multirow}
\usepackage{tabularx}

% Gráficos e Imagens
\usepackage{graphicx}
\usepackage{caption}
\usepackage{subcaption}

% Algoritmos e Listagens de Código
\usepackage{algorithm}
\usepackage{algpseudocode}
\usepackage{listings}
\usepackage{xcolor}

% Bibliografia e Citações (BibLaTeX com estilo IEEE)
\usepackage[backend=biber, style=numeric-comp, sorting=none]{biblatex}

% Hiperlinks (Sempre carregado por último)
\usepackage{hyperref}
\hypersetup{
    colorlinks=true,
    linkcolor=navyblue,
    citecolor=emerald,
    urlcolor=cobalt,
    pdftitle={Modelo Estruturado de Artigo Científico},
    pdfauthor={Nome do Autor}
}

% ==============================================================================
% DEFINIÇÃO DE CORES E ESTILOS
% ==============================================================================
\definecolor{navyblue}{RGB}{20, 50, 100}
\definecolor{emerald}{RGB}{0, 120, 80}
\definecolor{cobalt}{RGB}{0, 80, 180}
\definecolor{codegray}{rgb}{0.5,0.5,0.5}
\definecolor{codegreen}{rgb}{0,0.6,0}
\definecolor{codebackground}{rgb}{0.96,0.96,0.96}

% Configuração para listagens de código (Listings)
\lstset{
    backgroundcolor=\color{codebackground},
    commentstyle=\color{codegreen},
    keywordstyle=\color{navyblue}\bfseries,
    numberstyle=\tiny\color{codegray},
    stringstyle=\color{emerald},
    basicstyle=\ttfamily\small,
    breaklines=true,
    captionpos=b,
    keepspaces=true,
    numbers=left,
    numbersep=8pt,
    showspaces=false,
    showtabs=false
}

% ==============================================================================
% AMBIENTES MATEMÁTICOS (Teoremas, Definições, Lemas)
% ==============================================================================
\theoremstyle{definition}
\newtheorem{definition}{Definição}[section]
\newtheorem{theorem}{Teorema}[section]
\newtheorem{lemma}[theorem]{Lema}
\newtheorem{corollary}[theorem]{Corolário}

\theoremstyle{remark}
\newtheorem*{remark}{Observação}

% ==============================================================================
% BASE DE DADOS BIBLIOGRÁFICA (Exemplo Embutido)
% ==============================================================================
\begin{filecontents*}{\jobname.bib}
@article{shannon1948,
  author  = {Shannon, Claude E.},
  title   = {A Mathematical Theory of Communication},
  journal = {Bell System Technical Journal},
  volume  = {27},
  number  = {3},
  pages   = {379--423},
  year    = {1948}
}

@book{knuth1984,
  author    = {Knuth, Donald E.},
  title     = {The {\TeX}book},
  publisher = {Addison-Wesley},
  year      = {1984},
  address   = {Reading, Massachusetts}
}
\end{filecontents*}

\addbibresource{\jobname.bib}

% ==============================================================================
% METADADOS DO DOCUMENTO
% ==============================================================================
\title{\textbf{\Large Modelo Estruturado para Escrita Científica em {\LaTeX}}}
\author{
    \textbf{Autor Principal}\thanks{Contato: \href{mailto:autor@instituicao.org}{autor@instituicao.org}} \\
    \small Departamento de Ciência da Computação \\
    \small Universidade de Excelência Científica
    \and
    \textbf{Coautor da Pesquisa} \\
    \small Instituto de Física Aplicada \\
    \small Centro de Pesquisa Avançada
}
\date{\today}

% ==============================================================================
% INÍCIO DO DOCUMENTO
% ==============================================================================
\begin{document}

\maketitle

\begin{abstract}
Este documento apresenta um modelo abrangente e padronizado para a redação de artigos e relatórios científicos utilizando {\LaTeX}. São contempladas as principais funcionalidades necessárias na escrita acadêmica de alto nível: formulações matemáticas avançadas, estruturação de teoremas, tabelas no padrão \texttt{booktabs}, ambientes de código fonte, inclusão de algoritmos, figuras estruturadas e gerenciamento bibliográfico automatizado com \texttt{BibLaTeX}. O objetivo principal é servir como base reautilizável e pronta para submissões científicas.
\end{abstract}

\textbf{Palavras-chave:} Escrita Científica, {\LaTeX}, Modelo Estruturado, BibLaTeX, Matemática Formal.

\section{Introdução}
\label{sec:introducao}

A escrita científica exige rigor, clareza visual e reprodutibilidade. O sistema {\LaTeX} estabeleceu-se como o padrão de fato para a publicação em disciplinas exatas e tecnológicas devido ao seu controle preciso sobre tipografia e equações \cite{knuth1984}.

Neste modelo, cobrimos os blocos de construção essenciais que todo pesquisador utiliza ao compor um artigo ou relatório acadêmico:
\begin{itemize}
    \item Formatação formal de equações e ambientes de teoremas (Seção \ref{sec:matematica});
    \item Organização de dados quantitativos através de tabelas limpas (Seção \ref{sec:tabelas});
    \item Apresentação de pseudo-código e código de programação (Seção \ref{sec:algoritmos});
    \item Referência cruzada e citação bibliográfica rigorosa (Seção \ref{sec:referencias}).
\end{itemize}

\section{Fundamentação Teórica e Matemática}
\label{sec:matematica}

A manipulação de expressões matemáticas é uma das maiores forças do {\LaTeX}. Abaixo são demonstrados os ambientes matemáticos fundamentais disponibilizados pelos pacotes \texttt{amsmath} e \texttt{mathtools}.

\subsection{Equações Inline e Destacadas}

Variáveis e expressões simples podem ser escritas no corpo do texto, tais como $f(x) = x^2 + 2x + 1$ ou $e^{i\pi} + 1 = 0$. Para equações centrais ao desenvolvimento teórico, utiliza-se o ambiente numerado:

\begin{equation}
    \label{eq:fourier}
    \hat{f}(\xi) = \int_{-\infty}^{\infty} f(t) \, e^{-2\pi i t \xi} \, dt
\end{equation}

A Equação \ref{eq:fourier} representa a Transformada de Fourier contínua. Para sistemas de equações desalinhadas ou em múltiplas linhas, o ambiente \texttt{align} é recomendado:

\begin{align}
    \nabla \cdot \mathbf{E} &= \frac{\rho}{\varepsilon_0} \label{eq:maxwell1} \\
    \nabla \cdot \mathbf{B} &= 0 \label{eq:maxwell2} \\
    \nabla \times \mathbf{E} &= -\frac{\partial \mathbf{B}}{\partial t} \label{eq:maxwell3} \\
    \nabla \times \mathbf{B} &= \mu_0 \mathbf{J} + \mu_0 \varepsilon_0 \frac{\partial \mathbf{E}}{\partial t} \label{eq:maxwell4}
\end{align}

\subsection{Teoremas, Definições e Lemas}

Ambientes formais auxiliam na leitura e identificação de contribuições teóricas.

\begin{definition}[Entropia de Informação]
\label{def:entropia}
Seja $X$ uma variável aleatória discreta com conjunto de estados $\mathcal{X}$ e distribuição de probabilidade $P(X)$. A entropia $H(X)$ é dada por:
\begin{equation}
    H(X) = -\sum_{x \in \mathcal{X}} P(x) \log_2 P(x)
\end{equation}
\end{definition}

\begin{theorem}[Teorema da Amostragem de Nyquist-Shannon]
\label{thm:shannon}
Se um sinal contínuo $x(t)$ não contém frequências superiores a $B$ Hertz, ele pode ser completamente determinado a partir de suas amostras colhidas a uma taxa uniforme de $f_s > 2B$.
\end{theorem}

\begin{proof}
A demonstração decorre diretamente da análise de Fourier no domínio da frequência \cite{shannon1948}.
\end{proof}

\section{Apresentação de Dados: Tabelas e Figuras}
\label{sec:tabelas}

\subsection{Tabelas de Qualidade de Publicação}

Em documentos científicos impressos ou em PDF, deve-se evitar linhas verticais em tabelas. O pacote \texttt{booktabs} oferece os comandos \texttt{\textbackslash toprule}, \texttt{\textbackslash midrule} e \texttt{\textbackslash bottomrule}.

\begin{table}[h]
\centering
\caption{Comparação de desempenho entre algoritmos de ordenação.}
\label{tab:desempenho}
\begin{tabular}{l c c c}
\toprule
\textbf{Algoritmo} & \textbf{Complexidade (Méd)} & \textbf{Complexidade (Pior)} & \textbf{Memória} \\
\midrule
Quicksort   & $\mathcal{O}(n \log n)$ & $\mathcal{O}(n^2)$      & $\mathcal{O}(\log n)$ \\
Mergesort   & $\mathcal{O}(n \log n)$ & $\mathcal{O}(n \log n)$ & $\mathcal{O}(n)$ \\
Heapsort    & $\mathcal{O}(n \log n)$ & $\mathcal{O}(n \log n)$ & $\mathcal{O}(1)$ \\
\bottomrule
\end{tabular}
\end{table}

\subsection{Figuras e Subfiguras}

Para a inserção de elementos visuais, recomenda-se o uso de vetores (SVG convertidos para PDF) ou imagens de alta resolução (PNG/JPEG). O código a seguir exemplifica a estrutura formal (com uma caixa simulada para fins de demonstração).

\begin{figure}[h]
\centering
\framebox[0.6\textwidth][c]{\rule{0pt}{3cm} Representação Visual / Gráfico em PDF}
\caption{Diagrama conceitual do fluxo de processamento de dados.}
\label{fig:fluxo}
\end{figure}

\section{Algoritmos e Código-Fonte}
\label{sec:algoritmos}

O pacote \texttt{algpseudocode} permite descrever métodos procedimentais de forma estruturada.

\begin{algorithm}[h]
\caption{Busca Binária Iterativa}\label{alg:binarysearch}
\begin{algorithmic}[1]
\Procedure{BinarySearch}{$A, v, n$}
    \State $left \gets 1$
    \State $right \gets n$
    \While{$left \le right$}
        \State $mid \gets \lfloor (left + right) / 2 \rfloor$
        \If{$A[mid] == v$}
            \State \Return $mid$
        \ElsIf{$A[mid] < v$}
            \State $left \gets mid + 1$
        \Else
            \State $right \gets mid - 1$
        \EndIf
    \EndWhile
    \State \Return \textbf{null}
\EndProcedure
\end{algorithmic}
\end{algorithm}

Para trechos de código real, o pacote \texttt{listings} oferece realce de sintaxe personalizável:

\begin{lstlisting}[language=Python, caption={Exemplo de implementação em Python.}, label={lst:python}]
def calcular_entropia(probabilidades):
    import math
    entropia = 0.0
    for p in probabilidades:
        if p > 0:
            entropia -= p * math.log2(p)
    return entropia
\end{lstlisting}

\section{Citações e Referências Cruzadas}
\label{sec:referencias}

O gerenciamento de citações neste modelo utiliza \texttt{BibLaTeX}. É possível realizar citações diretas no texto utilizando o comando \texttt{\textbackslash cite}, por exemplo:
\begin{itemize}
    \item Trabalho pioneiro de Shannon sobre Teoria da Informação \cite{shannon1948};
    \item Referência essencial para tipografia técnica em {\LaTeX} \cite{knuth1984}.
\end{itemize}

Todas as referências citadas são compiladas automaticamente ao final do documento conforme a base de dados bibtex.

\section{Conclusão}
\label{sec:conclusao}

Este modelo fornece uma estrutura sólida e extensível para a produção de documentos científicos de alta qualidade. Com suporte nativo a equações complexas, tipografia rigorosa e gestão bibliográfica automatizada, o pesquisador pode focar no conteúdo científico mantendo a garantia de conformidade com os padrões de publicação internacional.

% ==============================================================================
% IMPRESSÃO DA BIBLIOGRAFIA
% ==============================================================================
\printbibliography[title={Referências Bibliográficas}]

\end{document}